% ─────────────────────────────────────────────────────────────────────────────
\section{Implementation, Tools, and Technologies}
\label{sec:implementation}
% ─────────────────────────────────────────────────────────────────────────────

The project was implemented by combining open-source vision,
multimodal reasoning, and edge-deployment techniques. The main technical challenge was
not inventing a new model from scratch, but selecting and integrating the right tools for the
problem. As the low-level design matured, the implementation focus expanded from model
selection alone to reproducible packaging, hardware-aware optimization, and explicit runtime
contracts for Jetson-class devices.

\subsection{Main Technologies Used}

\begin{table}[H]
  \centering
  \caption{Core tools and technologies used in the project}
  \begin{tabularx}{\textwidth}{@{}p{2.8cm}p{2.8cm}X@{}}
    \toprule
    Layer & Technology & Role in the project \\
    \midrule
    Detection & YOLOv8 & Fast object detection on thermal imagery with strong practical
    tooling and broad community support~\cite{yolov8}. \\
    Reasoning & Qwen2-VL 7B & Second-stage vision-language reasoning model used to
    interpret cropped thermal targets in more detail~\cite{qwen2vl}. \\
    Adaptation & LoRA & Parameter-efficient fine-tuning strategy for adapting the reasoning
    model without retraining all weights~\cite{lora}. \\
    Containerization & Docker / buildx & Reproducible deployment environment and cross-
    architecture image generation from workstation-class development hosts to ARM-based edge
    targets~\cite{dockerdocs}. \\
    Optimization runtime & TensorRT / Torch-TensorRT & Hardware-aware inference
    acceleration through engine compilation, mixed precision, and memory-conscious execution
    strategies~\cite{tensorrt,torchtrt}. \\
    Image processing & OpenCV & Region cropping, pre-processing, and data preparation
    utilities~\cite{opencv}. \\
    Deep learning runtime & PyTorch & Model development, loading, and experimentation
    framework~\cite{pytorch}. \\
    Deployment target & NVIDIA Jetson family & Edge-oriented hardware target for local,
    low-latency inference in a closed network~\cite{jetson}. \\
    Data interchange & JSON & Structured output format for alerts, logs, and downstream
    consumption. \\
    \bottomrule
  \end{tabularx}
\end{table}

\subsection{Detector Layer}

YOLOv8 was selected for the detector layer because the problem requires rapid scanning of an
entire scene before any deeper reasoning becomes possible. The model family is practical for
prototype work because it has mature training and inference tooling, strong existing
documentation, and published success on related thermal or aerial detection tasks.

The detector is responsible for:

\begin{itemize}[nosep]
  \item locating candidate targets in the full frame,
  \item producing bounding boxes and coarse class labels,
  \item acting as the gatekeeper for the more expensive reasoning stage.
\end{itemize}

This means the detector is not merely a preprocessing step; it is the part of the system that
protects runtime performance and gives the VLM a focused visual problem.

From a deployment perspective, the detector is also the component that benefits most directly
from TensorRT-style optimization because it is invoked continuously on the full frame and
therefore dominates the first stage of the runtime budget.

\subsection{Reasoning Layer}

Qwen2-VL was selected as the primary reasoning model because it provides strong open
multimodal performance, supports high-resolution inputs, and is well suited for prompt-based
description and inference tasks~\cite{qwen2vl}. In this project, the role of the model is not
to rediscover the object location but to examine the isolated target and produce a more useful
interpretation.

Typical reasoning goals include:

\begin{itemize}[nosep]
  \item narrowing ``vehicle'' to a more specific category such as pickup-type or light utility
    vehicle,
  \item inferring whether a strong localized thermal signature suggests recent engine activity,
  \item interpreting whether a human figure appears to be carrying a rifle-like object,
  \item converting low-level visual evidence into natural-language explanation.
\end{itemize}

Alternative open models such as LLaVA-NeXT and InternVL were also studied
as comparison points because they are strong open multimodal baselines with useful
reasoning capability~\cite{llavanext,internvl2024}. However, Qwen2-VL remained the most
aligned with the project's need for fine-grained visual reasoning under a relatively efficient
open-source setup.

\subsection{LoRA Fine-Tuning and Domain Adaptation}

One of the most important technical lessons of the project was that thermal adaptation is a
separate engineering problem from general multimodal reasoning. A VLM trained mostly on
ordinary web images does not automatically understand that thermal brightness represents
heat, not visible texture or color.

LoRA was therefore important as a strategy because it enables targeted adaptation of a large
model with lower computational cost than full retraining~\cite{lora}. Even when the final
prototype is evaluated primarily at system level, the project's design logic remained clear:
the analyst model benefits from thermal-domain adaptation, especially when asked to produce
object-specific explanations rather than generic captions. This emphasis on adaptation is also
consistent with broader thermal-domain transfer work, where modality mismatch is one of the
core causes of performance loss~\cite{sstn2021}.

\subsection{Containerized Deployment Workflow}

To make the implementation portable across workstation and Jetson environments, the low-
level deployment workflow was organized around containerization. This has several practical
advantages:

\begin{itemize}[nosep]
  \item model dependencies are frozen in a reproducible runtime image,
  \item Jetson-specific library versions can be aligned with the target JetPack or L4T stack,
  \item the same application pipeline can be developed on x86 hosts and prepared for ARM
    deployment through cross-architecture image builds.
\end{itemize}

The use of \texttt{docker buildx} is particularly valuable in this context because it bridges the
gap between high-resource development machines and ARM-based edge targets. Instead of
compiling every dependency directly on the embedded device, the build process can produce a
more controlled and repeatable deployment artifact in advance.

\subsection{TensorRT Optimization Strategy}

The implementation protocol treated acceleration as a structured workflow rather than a single
toggle. In edge deployment, optimization typically proceeds through three related steps:

\begin{enumerate}[nosep]
  \item graph-level simplification and layer fusion where supported,
  \item reduced-precision execution such as FP16 for lower latency and smaller memory use,
  \item engine-level runtime specialization for the target NVIDIA hardware.
\end{enumerate}

In this project, such optimizations are most naturally aligned with the detector stage and other
repeated tensor operations. For the reasoning stage, the analogous concerns are quantization,
reduced memory footprint, staged loading, and careful balancing between semantic depth and
runtime cost.

\subsection{Supporting Software Components}

OpenCV and PyTorch were used as the main supporting frameworks. OpenCV fits the bridge
layer of the architecture because cropping and preparing regions of interest is a direct image-
processing task. PyTorch fits the model layer because it is the de facto standard framework
for open detector and multimodal model experimentation.

JSON was used as the common output format because it offers a simple contract between
vision, reasoning, logging, and interface components. A structured format was necessary to
avoid producing outputs that were only human-readable but difficult to archive or integrate.

At implementation level, these supporting components also define memory boundaries. Raw
frames, normalized detector tensors, cropped ROIs, analysis objects, and final JSON records
must move between packages without ambiguity over format or ownership.

\subsection{Edge Optimization}

The project targeted edge-oriented operation rather than cloud dependence. This led to an
interest in quantization and model-size reduction strategies, especially for the reasoning
model. In principle, INT4-style quantization is attractive because it reduces memory pressure
and improves the feasibility of local deployment on Jetson-class hardware. This does not
eliminate the challenge of multimodal inference at the edge, but it makes the deployment goal
far more realistic.

Recent Jetson benchmarking work also shows that throughput, memory headroom, and energy
budget can vary substantially across embedded hardware classes, which is why edge AI design
cannot treat deployment as an afterthought~\cite{pham2023}. For this reason, the project
considered memory pressure and runtime continuity as first-class concerns alongside model
accuracy.

\subsection{Memory Management and Runtime Discipline}

Intermediate buffers can become as important as model weights in a
resource-constrained deployment. The implementation protocol therefore benefits from:

\begin{itemize}[nosep]
  \item explicit staging between frame ingestion, detection, cropping, and reasoning,
  \item avoiding unnecessary duplication of image tensors,
  \item conditional execution of the reasoning stage only when the detector produces targets,
  \item reduced-precision and engine-compiled execution where it provides a stable gain.
\end{itemize}

\subsection{Why These Technologies Matter}

Each tool solved a specific problem:

\begin{itemize}[nosep]
  \item YOLOv8 solved rapid thermal localization.
  \item Qwen2-VL solved semantic interpretation beyond labels.
  \item LoRA solved practical domain adaptation.
  \item OpenCV solved region isolation and pre-processing.
  \item PyTorch solved model experimentation and integration.
  \item Jetson-oriented optimization solved the closed-network edge-deployment goal.
\end{itemize}
